\documentclass[11pt,a4paper]{article}
\usepackage[utf8]{inputenc}
\usepackage[T1]{fontenc}
\usepackage[catalan]{babel}
\usepackage{lmodern}
\usepackage[sfdefault,lining]{FiraSans}
\usepackage[a4paper,top=2.3cm,bottom=2.2cm,left=1.8cm,right=1.8cm,headheight=1.3cm,headsep=0.4cm,footskip=1cm]{geometry}
\usepackage{amsmath,amssymb}
\usepackage{icomma}
\usepackage{graphicx}
\usepackage[table]{xcolor}
\usepackage{tikz}
\usepackage{fontawesome5}
\usepackage{ccicons}
\usepackage{fancyhdr}
\usepackage{lastpage}
\usepackage{tasks}
\usepackage{needspace}
\usepackage[most]{tcolorbox}
\usepackage{array}
\usepackage[hidelinks,pdftitle={Solucions de la fitxa d'exercicis de la unitat 2: Nombres reals. 3r d'ESO},pdfsubject={Matemàtiques 3r ESO},pdfkeywords={nombres reals, potències, radicals, notació científica, intervals}]{hyperref}

\definecolor{navy}{HTML}{1B2748}
\definecolor{accent}{HTML}{E8603A}
\definecolor{teal}{HTML}{1F9E90}
\definecolor{soft}{HTML}{EEF1F8}
\definecolor{muted}{HTML}{5D6685}

\setlength{\parindent}{0pt}
\setlength{\parskip}{3pt}
\settasks{label=\alph*), label-format=\bfseries\color{navy}, label-width=1.3em, item-indent=1.8em, column-sep=0.6em, after-item-skip=2pt, before-skip=3pt, after-skip=4pt}

% ---------- capçalera i peu ----------
\newcommand{\llicencia}{\href{https://creativecommons.org/licenses/by-nc-sa/4.0/deed.ca}{\ccbyncsa}\ \ Llicència Creative Commons Reconeixement-NoComercial-CompartirIgual 4.0 Internacional (CC BY-NC-SA 4.0)}
\pagestyle{fancy}
\fancyhf{}
\fancyhead[L]{\raisebox{-0.35\height}{\includegraphics[height=0.95cm]{logo-gva.png}}}
\fancyhead[C]{\small\color{muted}Solucions de la fitxa · Unitat 2: Nombres reals · 3r d'ESO}
\fancyhead[R]{\raisebox{-0.35\height}{\includegraphics[height=0.95cm]{logo-vermellar.png}}}
\fancyfoot[L]{\scriptsize\color{muted}\llicencia}
\fancyfoot[R]{\small\color{muted}\thepage}
\renewcommand{\headrulewidth}{0.4pt}
\renewcommand{\headrule}{\color{muted!40}\hrule height \headrulewidth}
\fancypagestyle{primera}{\fancyhf{}\renewcommand{\headrulewidth}{0pt}
  \fancyfoot[L]{\scriptsize\color{muted}\llicencia}
  \fancyfoot[R]{\small\color{muted}\thepage}}

% ---------- nivells i exercicis ----------
\newcommand{\estrelles}[1]{{\color{accent}\ifcase#1\or\faStar\faStar[regular]\faStar[regular]\or\faStar\faStar\faStar[regular]\or\faStar\faStar\faStar\fi}}
\newcounter{exer}
\newcommand{\ex}[1]{\par\needspace{4\baselineskip}\medskip\stepcounter{exer}\noindent{\color{navy}\bfseries\theexer.}\ \ {\small\estrelles{#1}}\ \ }

\newcommand{\seccio}[2]{\par\needspace{10\baselineskip}\bigskip
  \begin{tcolorbox}[colback=navy,colframe=navy,coltext=white,arc=2mm,boxrule=0pt,left=3mm,right=3mm,top=1.2mm,bottom=1.2mm]
  {\large\bfseries #1\quad #2}
  \end{tcolorbox}\vspace{-2pt}}

% ---------- recta real ----------
\newcommand{\rectabuida}{\begin{tikzpicture}[x=0.42cm,baseline=0]
  \draw[-latex,muted] (-6.7,0)--(6.8,0);
  \foreach \i in {-6,...,6}{\draw[muted] (\i,0.12)--(\i,-0.12) node[below=-1pt,font=\scriptsize,text=muted]{$\i$};}
  \end{tikzpicture}}
% \rectaint{x1}{x2}{estil esquerra}{estil dreta}: estil = ple, buit, inf
\newcommand{\extrem}[2]{\ifx#2p\fill[accent] (#1,0) circle (0.13);\else\ifx#2b\filldraw[fill=white,draw=accent,line width=1.2pt] (#1,0) circle (0.13);\fi\fi}
\newcommand{\rectaint}[4]{\begin{tikzpicture}[x=0.42cm,baseline=0]
  \draw[-latex,muted] (-6.7,0)--(6.8,0);
  \foreach \i in {-6,...,6}{\draw[muted] (\i,0.12)--(\i,-0.12) node[below=-1pt,font=\scriptsize,text=muted]{$\i$};}
  \draw[accent,line width=3pt] (#1,0)--(#2,0);
  \ifx#3i\fill[accent] (#1-0.25,0)--(#1+0.15,0.2)--(#1+0.15,-0.2)--cycle;\else\extrem{#1}{#3}\fi
  \ifx#4i\fill[accent] (#2+0.25,0)--(#2-0.15,0.2)--(#2-0.15,-0.2)--cycle;\else\extrem{#2}{#4}\fi
  \end{tikzpicture}}

\newcolumntype{C}[1]{>{\centering\arraybackslash}p{#1}}

% \solt: requadre per al resultat
\newcommand{\res}[1]{\boxed{#1}}
\newcommand{\rectaunio}{\begin{tikzpicture}[x=0.42cm,baseline=0]
  \draw[-latex,muted] (-6.7,0)--(6.8,0);
  \foreach \i in {-6,...,6}{\draw[muted] (\i,0.12)--(\i,-0.12) node[below=-1pt,font=\scriptsize,text=muted]{$\i$};}
  \draw[accent,line width=3pt] (-6.4,0)--(-1,0); \fill[accent] (-6.65,0)--(-6.25,0.2)--(-6.25,-0.2)--cycle; \fill[accent] (-1,0) circle (0.13);
  \draw[accent,line width=3pt] (4,0)--(6.4,0); \fill[accent] (6.65,0)--(6.25,0.2)--(6.25,-0.2)--cycle; \filldraw[fill=white,draw=accent,line width=1.2pt] (4,0) circle (0.13);
  \end{tikzpicture}}

% =====================================================================
% Capçalera FentMates de les fitxes: un color per a cada curs.
% Els colors són els de la pàgina inicial (LEVELS a assets/js/materials.js);
% el logo ve de _eines/logo/pdf/ (python3 _eines/logo/genera.py).
%
% Ús, al final del preàmbul de la unitat:
%   % =====================================================================
% Capçalera FentMates de les fitxes: un color per a cada curs.
% Els colors són els de la pàgina inicial (LEVELS a assets/js/materials.js);
% el logo ve de _eines/logo/pdf/ (python3 _eines/logo/genera.py).
%
% Ús, al final del preàmbul de la unitat:
%   % =====================================================================
% Capçalera FentMates de les fitxes: un color per a cada curs.
% Els colors són els de la pàgina inicial (LEVELS a assets/js/materials.js);
% el logo ve de _eines/logo/pdf/ (python3 _eines/logo/genera.py).
%
% Ús, al final del preàmbul de la unitat:
%   \input{../comu/fentmates.tex}
%   \fmcurs{1eso}
% i, al començament del document:
%   \capcalera{Matemàtiques}{1r ESO}{Fitxa d'exercicis · Unitat 2}{Divisibilitat}
% Opcional abans de \input: \def\fmneed{14} (espai mínim abans de cada secció).
% =====================================================================
\makeatletter
\def\fm@logodir{../../logo/pdf}
\providecommand{\fmneed}{10}
\def\fm@cursnom{}

% clau del curs, color i color del text damunt del color
\newcommand{\fm@curs}[3]{\@namedef{fm@color@#1}{#2}\@namedef{fm@text@#1}{#3}}
\fm@curs{1eso}{E6B84A}{navy}
\fm@curs{2eso}{3B82F6}{white}
\fm@curs{3eso}{FF7A50}{white}
\fm@curs{4eso-a}{A78BFA}{white}
\fm@curs{4eso-b}{6D28D9}{white}
\fm@curs{1batx-mi}{9B1C1C}{white}
\fm@curs{1batx-ccss}{EF6B6B}{white}
\fm@curs{2batx-mii}{1E6B47}{white}
\fm@curs{2batx-ccss}{2BB3A3}{white}

\newcommand{\fmcurs}[1]{%
  \@ifundefined{fm@color@#1}{\PackageError{fentmates}{Curs desconegut: #1}{}}{}%
  \edef\fm@tmp{\noexpand\definecolor{curs}{HTML}{\@nameuse{fm@color@#1}}%
               \noexpand\colorlet{curstext}{\@nameuse{fm@text@#1}}}\fm@tmp
  \colorlet{cursfosc}{curs!72!black}%
  \def\fm@logo{\fm@logodir/vertical-#1.pdf}%
  \def\fm@marca{\fm@logodir/marca-#1.pdf}}

% ---------- capçalera de la primera pàgina ----------
% #1 matèria (petit), #2 curs, #3 tipus de document i unitat, #4 títol
\newcommand{\capcalera}[4]{%
  \gdef\fm@cursnom{#2}%
  \par\noindent
  \begin{tcolorbox}[colback=curs,colframe=curs,boxrule=0pt,arc=2.2mm,width=3.5cm,height=2.15cm,
      valign=center,left=3mm,right=2mm,top=1mm,bottom=1mm,nobeforeafter,box align=center]
    \color{curstext}{\footnotesize #1}\par\vspace{1pt}{\fontsize{19}{21}\selectfont\bfseries #2\par}%
  \end{tcolorbox}\hspace{5mm}%
  \begin{minipage}[c]{\dimexpr\linewidth-3.5cm-5mm-2.9cm\relax}
    {\color{cursfosc}\small\bfseries\MakeUppercase{#3}}\par\vspace{2pt}
    {\color{navy}\fontsize{20}{23}\selectfont\bfseries #4\par}
  \end{minipage}\hfill
  \begin{minipage}[c]{2.6cm}\raggedleft\includegraphics[width=2.6cm]{\fm@logo}\end{minipage}\par
  \vspace{2pt}}

% ---------- barres de secció del color del curs ----------
\renewcommand{\seccio}[2]{\par\needspace{\fmneed\baselineskip}\bigskip
  \begin{tcolorbox}[colback=curs,colframe=curs,coltext=curstext,arc=2mm,boxrule=0pt,left=3mm,right=3mm,top=1.2mm,bottom=1.2mm]
  {\large\bfseries #1\quad #2}
  \end{tcolorbox}\vspace{-2pt}}

% ---------- capçalera i peu de les pàgines ----------
% primera pàgina: logos institucionals al peu
\fancypagestyle{primera}{\fancyhf{}\renewcommand{\headrulewidth}{0pt}
  \fancyfoot[L]{\raisebox{-0.32\height}{\includegraphics[height=0.78cm]{logo-gva.png}}\hspace{5mm}%
    \raisebox{-0.32\height}{\includegraphics[height=0.78cm]{logo-vermellar.png}}}
  \fancyfoot[R]{\scriptsize\color{muted}\href{https://creativecommons.org/licenses/by-nc-sa/4.0/deed.ca}{\ccbyncsa}\ \ CC BY-NC-SA 4.0 · fentmates.org\qquad\small\thepage}}
% la resta: el símbol a l'esquerra i el curs a la dreta, amb el títol al centre com abans
\fancyhead[L]{\raisebox{-0.3\height}{\includegraphics[height=0.62cm]{\fm@marca}}}
\fancyhead[R]{\small\bfseries\color{cursfosc}\fm@cursnom}
\renewcommand{\headrule}{\color{curs!45}\hrule height \headrulewidth}
\makeatother

%   \fmcurs{1eso}
% i, al començament del document:
%   \capcalera{Matemàtiques}{1r ESO}{Fitxa d'exercicis · Unitat 2}{Divisibilitat}
% Opcional abans de \input: \def\fmneed{14} (espai mínim abans de cada secció).
% =====================================================================
\makeatletter
\def\fm@logodir{../../logo/pdf}
\providecommand{\fmneed}{10}
\def\fm@cursnom{}

% clau del curs, color i color del text damunt del color
\newcommand{\fm@curs}[3]{\@namedef{fm@color@#1}{#2}\@namedef{fm@text@#1}{#3}}
\fm@curs{1eso}{E6B84A}{navy}
\fm@curs{2eso}{3B82F6}{white}
\fm@curs{3eso}{FF7A50}{white}
\fm@curs{4eso-a}{A78BFA}{white}
\fm@curs{4eso-b}{6D28D9}{white}
\fm@curs{1batx-mi}{9B1C1C}{white}
\fm@curs{1batx-ccss}{EF6B6B}{white}
\fm@curs{2batx-mii}{1E6B47}{white}
\fm@curs{2batx-ccss}{2BB3A3}{white}

\newcommand{\fmcurs}[1]{%
  \@ifundefined{fm@color@#1}{\PackageError{fentmates}{Curs desconegut: #1}{}}{}%
  \edef\fm@tmp{\noexpand\definecolor{curs}{HTML}{\@nameuse{fm@color@#1}}%
               \noexpand\colorlet{curstext}{\@nameuse{fm@text@#1}}}\fm@tmp
  \colorlet{cursfosc}{curs!72!black}%
  \def\fm@logo{\fm@logodir/vertical-#1.pdf}%
  \def\fm@marca{\fm@logodir/marca-#1.pdf}}

% ---------- capçalera de la primera pàgina ----------
% #1 matèria (petit), #2 curs, #3 tipus de document i unitat, #4 títol
\newcommand{\capcalera}[4]{%
  \gdef\fm@cursnom{#2}%
  \par\noindent
  \begin{tcolorbox}[colback=curs,colframe=curs,boxrule=0pt,arc=2.2mm,width=3.5cm,height=2.15cm,
      valign=center,left=3mm,right=2mm,top=1mm,bottom=1mm,nobeforeafter,box align=center]
    \color{curstext}{\footnotesize #1}\par\vspace{1pt}{\fontsize{19}{21}\selectfont\bfseries #2\par}%
  \end{tcolorbox}\hspace{5mm}%
  \begin{minipage}[c]{\dimexpr\linewidth-3.5cm-5mm-2.9cm\relax}
    {\color{cursfosc}\small\bfseries\MakeUppercase{#3}}\par\vspace{2pt}
    {\color{navy}\fontsize{20}{23}\selectfont\bfseries #4\par}
  \end{minipage}\hfill
  \begin{minipage}[c]{2.6cm}\raggedleft\includegraphics[width=2.6cm]{\fm@logo}\end{minipage}\par
  \vspace{2pt}}

% ---------- barres de secció del color del curs ----------
\renewcommand{\seccio}[2]{\par\needspace{\fmneed\baselineskip}\bigskip
  \begin{tcolorbox}[colback=curs,colframe=curs,coltext=curstext,arc=2mm,boxrule=0pt,left=3mm,right=3mm,top=1.2mm,bottom=1.2mm]
  {\large\bfseries #1\quad #2}
  \end{tcolorbox}\vspace{-2pt}}

% ---------- capçalera i peu de les pàgines ----------
% primera pàgina: logos institucionals al peu
\fancypagestyle{primera}{\fancyhf{}\renewcommand{\headrulewidth}{0pt}
  \fancyfoot[L]{\raisebox{-0.32\height}{\includegraphics[height=0.78cm]{logo-gva.png}}\hspace{5mm}%
    \raisebox{-0.32\height}{\includegraphics[height=0.78cm]{logo-vermellar.png}}}
  \fancyfoot[R]{\scriptsize\color{muted}\href{https://creativecommons.org/licenses/by-nc-sa/4.0/deed.ca}{\ccbyncsa}\ \ CC BY-NC-SA 4.0 · fentmates.org\qquad\small\thepage}}
% la resta: el símbol a l'esquerra i el curs a la dreta, amb el títol al centre com abans
\fancyhead[L]{\raisebox{-0.3\height}{\includegraphics[height=0.62cm]{\fm@marca}}}
\fancyhead[R]{\small\bfseries\color{cursfosc}\fm@cursnom}
\renewcommand{\headrule}{\color{curs!45}\hrule height \headrulewidth}
\makeatother

%   \fmcurs{1eso}
% i, al començament del document:
%   \capcalera{Matemàtiques}{1r ESO}{Fitxa d'exercicis · Unitat 2}{Divisibilitat}
% Opcional abans de \input: \def\fmneed{14} (espai mínim abans de cada secció).
% =====================================================================
\makeatletter
\def\fm@logodir{../../logo/pdf}
\providecommand{\fmneed}{10}
\def\fm@cursnom{}

% clau del curs, color i color del text damunt del color
\newcommand{\fm@curs}[3]{\@namedef{fm@color@#1}{#2}\@namedef{fm@text@#1}{#3}}
\fm@curs{1eso}{E6B84A}{navy}
\fm@curs{2eso}{3B82F6}{white}
\fm@curs{3eso}{FF7A50}{white}
\fm@curs{4eso-a}{A78BFA}{white}
\fm@curs{4eso-b}{6D28D9}{white}
\fm@curs{1batx-mi}{9B1C1C}{white}
\fm@curs{1batx-ccss}{EF6B6B}{white}
\fm@curs{2batx-mii}{1E6B47}{white}
\fm@curs{2batx-ccss}{2BB3A3}{white}

\newcommand{\fmcurs}[1]{%
  \@ifundefined{fm@color@#1}{\PackageError{fentmates}{Curs desconegut: #1}{}}{}%
  \edef\fm@tmp{\noexpand\definecolor{curs}{HTML}{\@nameuse{fm@color@#1}}%
               \noexpand\colorlet{curstext}{\@nameuse{fm@text@#1}}}\fm@tmp
  \colorlet{cursfosc}{curs!72!black}%
  \def\fm@logo{\fm@logodir/vertical-#1.pdf}%
  \def\fm@marca{\fm@logodir/marca-#1.pdf}}

% ---------- capçalera de la primera pàgina ----------
% #1 matèria (petit), #2 curs, #3 tipus de document i unitat, #4 títol
\newcommand{\capcalera}[4]{%
  \gdef\fm@cursnom{#2}%
  \par\noindent
  \begin{tcolorbox}[colback=curs,colframe=curs,boxrule=0pt,arc=2.2mm,width=3.5cm,height=2.15cm,
      valign=center,left=3mm,right=2mm,top=1mm,bottom=1mm,nobeforeafter,box align=center]
    \color{curstext}{\footnotesize #1}\par\vspace{1pt}{\fontsize{19}{21}\selectfont\bfseries #2\par}%
  \end{tcolorbox}\hspace{5mm}%
  \begin{minipage}[c]{\dimexpr\linewidth-3.5cm-5mm-2.9cm\relax}
    {\color{cursfosc}\small\bfseries\MakeUppercase{#3}}\par\vspace{2pt}
    {\color{navy}\fontsize{20}{23}\selectfont\bfseries #4\par}
  \end{minipage}\hfill
  \begin{minipage}[c]{2.6cm}\raggedleft\includegraphics[width=2.6cm]{\fm@logo}\end{minipage}\par
  \vspace{2pt}}

% ---------- barres de secció del color del curs ----------
\renewcommand{\seccio}[2]{\par\needspace{\fmneed\baselineskip}\bigskip
  \begin{tcolorbox}[colback=curs,colframe=curs,coltext=curstext,arc=2mm,boxrule=0pt,left=3mm,right=3mm,top=1.2mm,bottom=1.2mm]
  {\large\bfseries #1\quad #2}
  \end{tcolorbox}\vspace{-2pt}}

% ---------- capçalera i peu de les pàgines ----------
% primera pàgina: logos institucionals al peu
\fancypagestyle{primera}{\fancyhf{}\renewcommand{\headrulewidth}{0pt}
  \fancyfoot[L]{\raisebox{-0.32\height}{\includegraphics[height=0.78cm]{logo-gva.png}}\hspace{5mm}%
    \raisebox{-0.32\height}{\includegraphics[height=0.78cm]{logo-vermellar.png}}}
  \fancyfoot[R]{\scriptsize\color{muted}\href{https://creativecommons.org/licenses/by-nc-sa/4.0/deed.ca}{\ccbyncsa}\ \ CC BY-NC-SA 4.0 · fentmates.org\qquad\small\thepage}}
% la resta: el símbol a l'esquerra i el curs a la dreta, amb el títol al centre com abans
\fancyhead[L]{\raisebox{-0.3\height}{\includegraphics[height=0.62cm]{\fm@marca}}}
\fancyhead[R]{\small\bfseries\color{cursfosc}\fm@cursnom}
\renewcommand{\headrule}{\color{curs!45}\hrule height \headrulewidth}
\makeatother

\fmcurs{3eso}

\begin{document}
\thispagestyle{primera}

% ---------- encapçalament ----------
\capcalera{Matemàtiques}{3r ESO}{Solucions de la fitxa · Unitat 2}{Nombres reals}
\vspace{8pt}

\begin{tcolorbox}[colback=soft,colframe=soft,arc=2mm,left=3mm,right=3mm,top=1.5mm,bottom=1.5mm]
\small Intenta resoldre cada exercici abans de mirar-ne la solució. Si no t'ix el mateix resultat, repassa els passos o consulta els apunts interactius de la unitat. Quan hi ha diverses respostes possibles, se'n mostra una d'exemple.
\end{tcolorbox}

% =====================================================
\seccio{1}{Potències d'exponent natural}

\ex{1}
\begin{tasks}(4)
  \task $5^{3}=125$
  \task $(-2)^{5}=-32$
  \task $\left(\frac23\right)^{4}=\frac{16}{81}$
  \task $(-1)^{7}=-1$
\end{tasks}

\ex{1}
\begin{tasks}(4)
  \task $64$
  \task $-27$
  \task $81$
  \task $-81$
  \task $1$
  \task $-\frac18$
  \task $0$
  \task $10\,000\,000$
\end{tasks}

\ex{2}
\begin{tasks}(2)
  \task Positiu: base negativa i exponent parell.
  \task Negatiu: base negativa i exponent imparell.
  \task Negatiu: $(-2)^{6}=64$, i el signe de davant el canvia: $-64$.
  \task Positiu: $(-4)^{3}=-64$, i $-(-64)=64$.
\end{tasks}

\ex{2}
\begin{tasks}(2)
  \task $-4+4-(-8)=\res{8}$
  \task $-1+1-1=\res{-1}$
  \task $9-8+1=\res{2}$
  \task $\frac49\cdot\frac{27}{8}=\frac{108}{72}=\res{\frac32}$
\end{tasks}

\ex{3}
\begin{tasks}(4)
  \task $x=6$
  \task $x=3$
  \task $x=3$\ \ o\ \ $x=-3$
  \task $x=3$
\end{tasks}

% =====================================================
\seccio{2}{Propietats de les potències}

\ex{1}
\begin{tasks}(6)
  \task $3^{9}$
  \task $7^{4}$
  \task $2^{15}$
  \task $5^{8}$
  \task $a^{5}$
  \task $(-2)^{6}=2^{6}$
\end{tasks}

\ex{1}
\begin{tasks}(4)
  \task $10^{5}=100\,000$
  \task $3^{4}=81$
  \task $6^{3}=216$
  \task $1^{2}=1$
\end{tasks}

\ex{2}
\begin{tasks}(4)
  \task $\dfrac{2^{10}}{2^{8}}=2^{2}=\res{4}$
  \task $\dfrac{3^{11}}{3^{9}}=3^{2}=\res{9}$
  \task $\dfrac{5^{9}}{5^{8}}=\res{5}$
  \task $\dfrac{6^{8}}{6^{6}}=6^{2}=\res{36}$
\end{tasks}

\ex{2}
\begin{tasks}(2)
  \task $\dfrac{2^{6}\cdot2^{6}}{2^{8}}=2^{4}=\res{16}$
  \task $\dfrac{3^{4}\cdot3^{3}}{3^{6}}=\res{3}$
  \task $\dfrac{5^{6}}{5^{2}\cdot5^{3}}=\res{5}$
  \task $\dfrac{2^{4}\cdot3^{4}}{2^{3}\cdot3^{2}}=2\cdot3^{2}=\res{18}$
\end{tasks}

\ex{3}
\begin{tasks}(2)
  \task $\dfrac{2^{6}\cdot3^{3}\cdot2^{2}\cdot3^{4}}{2^{6}\cdot3^{6}}=\dfrac{2^{8}\cdot3^{7}}{2^{6}\cdot3^{6}}=2^{2}\cdot3=\res{12}$
  \task $\dfrac{3^{2}\cdot5^{2}\cdot2^{3}\cdot5^{3}}{2^{2}\cdot3^{2}\cdot5^{4}}=\dfrac{2^{3}\cdot3^{2}\cdot5^{5}}{2^{2}\cdot3^{2}\cdot5^{4}}=2\cdot5=\res{10}$
  \task $\dfrac{2^{6}\cdot5^{3}\cdot3^{4}}{2^{2}\cdot3^{2}\cdot5^{2}\cdot2^{4}\cdot3}=\dfrac{2^{6}\cdot3^{4}\cdot5^{3}}{2^{6}\cdot3^{3}\cdot5^{2}}=3\cdot5=\res{15}$
  \task $\dfrac{2^{2}\cdot7^{2}\cdot3^{3}\cdot7^{3}}{7^{4}\cdot2^{2}\cdot3^{2}}=\dfrac{2^{2}\cdot3^{3}\cdot7^{5}}{2^{2}\cdot3^{2}\cdot7^{4}}=3\cdot7=\res{21}$
\end{tasks}

\ex{3}
\begin{tasks}(2)
  \task Vertadera: $8+8=16=2^{4}$.
  \task Falsa: $3^{2}\cdot3^{3}=3^{5}$.
  \task Falsa: $25\neq4+9=13$.
  \task Vertadera: $(5\cdot2)^{3}=10^{3}$.
  \task Vertadera: $2\cdot2^{10}=2^{11}$.
  \task Falsa: $4^{5}:4^{5}=4^{0}=1$.
\end{tasks}

% =====================================================
\seccio{3}{Potències d'exponent enter negatiu}

\ex{1}
\begin{tasks}(8)
  \task $\frac19$
  \task $-\frac18$
  \task $\frac{1}{1000}$
  \task $16$
  \task $\frac{27}{8}$
  \task $\frac{25}{9}$
  \task $1$
  \task $-\frac{1}{16}$
\end{tasks}

\ex{1}
\begin{tasks}(5)
  \task $5^{-4}$
  \task $3^{-3}$
  \task $2^{-4}$\ \ (o $4^{-2}$)
  \task $10^{-3}$
  \task $7^{-2}$
\end{tasks}

\ex{2}
\begin{tasks}(6)
  \task $5^{4}$
  \task $2^{-7}$
  \task $3^{6}$
  \task $7^{2}$
  \task $4^{2}$
  \task $10^{-4}$
\end{tasks}

\ex{2}
\begin{tasks}(4)
  \task $\frac12+\frac14=\res{\frac34}$
  \task $\frac13-\frac19=\res{\frac29}$
  \task $4+3=\res{7}$
  \task $1-\frac15=\res{\frac45}$
\end{tasks}

\ex{3}
\begin{tasks}(2)
  \task $\dfrac{2^{-6}\cdot2^{6}}{2^{-5}\cdot2^{4}}=\dfrac{2^{0}}{2^{-1}}=\res{2}$
  \task $\dfrac{3^{-2}\cdot3^{6}}{3^{-2}\cdot3^{4}}=\dfrac{3^{4}}{3^{2}}=\res{9}$
  \task $\dfrac{2^{-2}\cdot3^{-2}\cdot2^{6}\cdot3^{3}}{2^{-1}\cdot3^{-2}\cdot2^{5}}=\dfrac{2^{4}\cdot3}{2^{4}\cdot3^{-2}}=3^{3}=\res{27}$
  \task $\dfrac{2^{-2}\cdot5^{-2}\cdot5^{6}}{5^{4}\cdot2^{-3}}=\dfrac{2^{-2}\cdot5^{4}}{2^{-3}\cdot5^{4}}=\res{2}$
  \task $\dfrac{3^{2}}{2^{2}}\cdot\dfrac{2^{2}}{3^{2}}=\res{1}$
  \task $\left[2^{3}\cdot2^{-4}\right]^{-1}=\left(2^{-1}\right)^{-1}=\res{2}$
\end{tasks}

% =====================================================
\seccio{4}{Notació científica}

\ex{1}
\begin{tasks}(6)
  \task $4,5\cdot10^{7}$
  \task $7,2\cdot10^{-4}$
  \task $9,08\cdot10^{5}$
  \task $3\cdot10^{-7}$
  \task $-1,25\cdot10^{4}$
  \task $5\cdot10^{-2}$
\end{tasks}

\ex{1}
\begin{tasks}(5)
  \task $3\,200\,000$
  \task $0,005$
  \task $1\,070\,000\,000$
  \task $0,000\,008\,45$
  \task $-61\,000$
\end{tasks}

\ex{2} L'apartat a) ja està en notació científica. La resta, corregits:
\begin{tasks}[start=2](5)
  \task $3,4\cdot10^{7}$
  \task $8\cdot10^{-4}$
  \task $1,05\cdot10^{-1}$
  \task $1,25\cdot10^{-6}$
  \task $4,2\cdot10^{6}$
\end{tasks}

\ex{2}
\[9\cdot10^{-6}\ <\ 5\cdot10^{-5}\ <\ 1,2\cdot10^{-4}\ <\ 3,1\cdot10^{-4}\ <\ 2,9\cdot10^{-3}\]

\ex{3}
\begin{tasks}(3)
  \task $3\cdot10^{5}$ km/s
  \task $7,5\cdot10^{-6}$ m
  \task $8\cdot10^{9}$ persones
  \task $1,67\cdot10^{-24}$ g
  \task $9,46\cdot10^{12}$ km
\end{tasks}

% =====================================================
\seccio{5}{Radicals. Extracció i introducció de factors}

\ex{1}
\begin{tasks}(4)
  \task $9$
  \task $-5$
  \task $2$
  \task No existeix en $\mathbb{R}$
  \task $-2$
  \task $\frac57$
  \task $\frac23$
  \task $1$
\end{tasks}

\ex{1}
\begin{tasks}(6)
  \task $2\sqrt{3}$
  \task $3\sqrt{2}$
  \task $3\sqrt{5}$
  \task $7\sqrt{2}$
  \task $2\sqrt[3]{2}$
  \task $3\sqrt[3]{3}$
\end{tasks}

\ex{2}
\begin{tasks}(3)
  \task $\sqrt{2^{3}\cdot5^{2}}=10\sqrt{2}$
  \task $2^{2}\cdot3\sqrt{2\cdot3}=12\sqrt{6}$
  \task $2^{2}\cdot5\sqrt[3]{2\cdot5}=20\sqrt[3]{10}$
  \task $3^{2}\sqrt[4]{3}=9\sqrt[4]{3}$
  \task $a^{3}b^{2}\sqrt{a}$
  \task $x\,y^{3}\sqrt[3]{x^{2}}$
\end{tasks}

\ex{3}
\begin{tasks}(4)
  \task $\dfrac{5\sqrt{2}}{9}$
  \task $2a\sqrt{3a}$
  \task $-3\sqrt[3]{2}$
  \task $2ab\sqrt[4]{2a^{2}b}$
  \task $\dfrac{6\sqrt{2}}{5}$
  \task $-2a\sqrt[3]{2a}$
  \task $3a^{2}b\sqrt{2a}$
  \task $3x^{2}y\sqrt[3]{3x}$
\end{tasks}

\ex{1}
\begin{tasks}(6)
  \task $\sqrt{20}$
  \task $\sqrt{27}$
  \task $\sqrt{50}$
  \task $\sqrt[3]{24}$
  \task $\sqrt[4]{162}$
  \task $\sqrt{a^{2}b}$
\end{tasks}
\begin{tasks}[resume](2)
  \task $\sqrt{\left(2^{3}\right)^{2}\cdot2}=\sqrt{2^{6}\cdot2}=\res{\sqrt{2^{7}}}$
  \task $\sqrt{\left(3^{2}\right)^{2}\cdot5}=\sqrt{3^{4}\cdot5}=\res{\sqrt{405}}$
  \task $\sqrt[3]{\left(2^{2}\right)^{3}\cdot2}=\sqrt[3]{2^{6}\cdot2}=\res{\sqrt[3]{2^{7}}}$
  \task $\sqrt[3]{\left(5^{2}\right)^{3}\cdot5^{2}}=\sqrt[3]{5^{6}\cdot5^{2}}=\res{\sqrt[3]{5^{8}}}$
  \task $\sqrt{\left(2^{3}\right)^{2}\cdot3^{2}\cdot3}=\sqrt{2^{6}\cdot3^{3}}=\res{\sqrt{1728}}$
  \task $\sqrt[4]{\left(a^{2}\right)^{4}\cdot a^{3}}=\sqrt[4]{a^{8}\cdot a^{3}}=\res{\sqrt[4]{a^{11}}}$
\end{tasks}

\ex{2} $3\sqrt5=\sqrt{45}$,\ \ $4\sqrt3=\sqrt{48}$,\ \ $2\sqrt{11}=\sqrt{44}$,\ \ $5\sqrt2=\sqrt{50}$. Per tant:
\[2\sqrt{11}\ <\ 3\sqrt{5}\ <\ 4\sqrt{3}\ <\ 5\sqrt{2}\]

% =====================================================
\seccio{6}{Operacions amb radicals}

\ex{1}
\begin{tasks}(4)
  \task $5\sqrt{5}$
  \task $4\sqrt{3}$
  \task $5\sqrt[3]{6}$
  \task $-2\sqrt{2}$
\end{tasks}

\ex{2}
\begin{tasks}(2)
  \task $2\sqrt2+3\sqrt2=\res{5\sqrt2}$
  \task $2\sqrt3+5\sqrt3-3\sqrt3=\res{4\sqrt3}$
  \task $4\sqrt5-3\sqrt5+4\sqrt5=\res{5\sqrt5}$
  \task $5\sqrt2-8\sqrt2+3\sqrt2=\res{0}$
  \task $2\sqrt[3]{2}+3\sqrt[3]{2}-\sqrt[3]{2}=\res{4\sqrt[3]{2}}$
  \task $2\sqrt7+3\sqrt7-\sqrt7=\res{4\sqrt7}$
\end{tasks}

\ex{1}
\begin{tasks}(6)
  \task $\sqrt{36}=6$
  \task $\sqrt{100}=10$
  \task $\sqrt{12}=2\sqrt{3}$
  \task $\sqrt[3]{8}=2$
  \task $\sqrt{16}=4$
  \task $\sqrt[3]{27}=3$
\end{tasks}

\ex{2}
\begin{tasks}(3)
  \task $10\sqrt{18}=\res{30\sqrt{2}}$
  \task $\sqrt{150}=\res{5\sqrt{6}}$
  \task $9\cdot2=\res{18}$
  \task $3\sqrt{5}$
  \task $\sqrt{3^{5}}=\res{9\sqrt{3}}$
  \task $\sqrt{36}=\res{6}$
\end{tasks}

\ex{3}
\begin{tasks}(2)
  \task $\sqrt{16}+\sqrt{36}=4+6=\res{10}$
  \task $2\cdot3-\sqrt{36}=6-6=\res{0}$
  \task $\sqrt{12}+2=\res{2\sqrt{3}+2}$
  \task $2\sqrt3+3\sqrt3-\sqrt{18}=\res{5\sqrt{3}-3\sqrt{2}}$
  \task $\dfrac{5\sqrt2-2\sqrt2}{\sqrt2}=\dfrac{3\sqrt2}{\sqrt2}=\res{3}$
\end{tasks}

% =====================================================
\seccio{7}{Intervals: definició i tipus}

\ex{1}\par\smallskip
{\renewcommand{\arraystretch}{1.45}\centering
\begin{tabular}{|C{3.6cm}|C{4.2cm}|C{6.2cm}|}
\hline
\rowcolor{soft}\textbf{Interval} & \textbf{Desigualtat} & \textbf{Tipus d'interval} \\ \hline
$(-3,\ 5)$ & $-3<x<5$ & Obert \\ \hline
$[-1,\ 4]$ & $-1\le x\le 4$ & Tancat \\ \hline
$(2,\ 6]$ & $2<x\le 6$ & Semiobert (obert per l'esquerra) \\ \hline
$[-2,\ +\infty)$ & $x\ge -2$ & Semirecta tancada \\ \hline
$(-\infty,\ 0)$ & $x<0$ & Semirecta oberta \\ \hline
$[-5,\ -1)$ & $-5\le x<-1$ & Semiobert (obert per la dreta) \\ \hline
\end{tabular}\par}

\ex{1}
\begin{tasks}(2)
  \task $(-4,\ +\infty)$;\ \ $x>-4$
  \task $(-\infty,\ 7]$;\ \ $x\le 7$
  \task $[0,\ 10]$;\ \ $0\le x\le 10$
  \task $(-2,\ 3]$;\ \ $-2<x\le 3$
\end{tasks}

\ex{2}
\begin{tasks}(3)
  \task No (el 3 no hi pertany)
  \task Sí
  \task Sí
  \task Sí
  \task No ($-7<-6$)
  \task Sí ($\sqrt2\approx1,41$)
\end{tasks}

\ex{3}
\begin{tasks}(1)
  \task Per exemple, $1,1$;\ \ $1,5$\ \ i\ \ $1,99$. L'interval conté infinits nombres reals.
  \task Per exemple, $(-1,\ 6)$ o $[0,\ 5]$.
  \task En $[-3,\ 4)$ és el $3$ (el 4 no hi pertany). En $(-3,\ 4]$ és el $4$.
  \task $(-1,\ 5)$, perquè $2-3=-1$ i $2+3=5$.
\end{tasks}

% =====================================================
\needspace{18\baselineskip}
\seccio{8}{Representació gràfica d'intervals}

\ex{1}\par\smallskip
\begin{tabular}{@{}l@{\ \ }l@{\qquad}l@{\ \ }l@{}}
a) $[-2,\ 3]$ & \rectaint{-2}{3}{p}{p} & b) $(0,\ 4)$ & \rectaint{0}{4}{b}{b} \\[16pt]
c) $(-\infty,\ 1]$ & \rectaint{-6.4}{1}{i}{p} & d) $(-3,\ +\infty)$ & \rectaint{-3}{6.4}{b}{i} \\[16pt]
e) $[-4,\ -1)$ & \rectaint{-4}{-1}{p}{b} & f) $x\le 2$ & \rectaint{-6.4}{2}{i}{p}
\end{tabular}

\ex{2}
\begin{tasks}(2)
  \task $(-1,\ 4]$;\ \ $-1<x\le 4$
  \task $[2,\ +\infty)$;\ \ $x\ge 2$
  \task $(-\infty,\ -3)$;\ \ $x<-3$
  \task $[-4,\ 0]$;\ \ $-4\le x\le 0$
\end{tasks}

\ex{3}\par\smallskip
\begin{tabular}{@{}l@{\ \ }l@{}}
a) $(-1,\ 5]$ & \rectaint{-1}{5}{b}{p} \\[16pt]
b) $[-3,\ 2)$ & \rectaint{-3}{2}{p}{b} \\[16pt]
c) $(-\infty,\ -1]\cup(4,\ +\infty)$ & \rectaunio
\end{tabular}

% =====================================================
\seccio{9}{Unió i intersecció d'intervals}

\ex{1}\par\smallskip
\begin{tabular}{@{}l@{\ \ }l@{}}
$A$ & \rectaint{-2}{4}{p}{p} \\[14pt]
$B$ & \rectaint{1}{6}{b}{b}
\end{tabular}\par\smallskip
$A\cup B=\res{[-2,\ 6)}\qquad A\cap B=\res{(1,\ 4]}$

\ex{2}
\begin{tasks}(2)
  \task $A\cup B=\mathbb{R}$;\ \ $A\cap B=[0,\ 3)$
  \task $A\cup B=[-5,\ 2)$;\ \ $A\cap B=(-3,\ -1]$
  \task $A\cup B=(-\infty,\ -2]\cup(1,\ +\infty)$;\ \ $A\cap B=\varnothing$
  \task $A\cup B=[0,\ 8]$;\ \ $A\cap B=\varnothing$
  \task $A\cup B=(-4,\ 4)$;\ \ $A\cap B=[-1,\ 1]$
\end{tasks}

\ex{3}
\begin{tasks}(1)
  \task Per exemple, $A=[0,\ 3]$ i $B=[2,\ 5]$.
  \task $a=1$.
  \task $A\cap B=(0,\ 6]$, i per tant $A\cap B\cap C=(0,\ 3]$\ \ i\ \ $(A\cap B)\cup C=(-\infty,\ 6]$.
  \task $[-1,\ 7)=[-1,\ +\infty)\cap(-\infty,\ 7)$.
\end{tasks}

% =====================================================
\seccio{10}{Classificació dels nombres reals}

\ex{1}\par\smallskip
{\renewcommand{\arraystretch}{1.3}\centering
\begin{tabular}{|C{3.6cm}|C{1.3cm}|C{1.3cm}|C{1.3cm}|C{1.3cm}|C{1.3cm}|}
\hline
\rowcolor{soft}\textbf{Nombre} & $\mathbb{N}$ & $\mathbb{Z}$ & $\mathbb{Q}$ & $\mathbb{I}$ & $\mathbb{R}$ \\ \hline
$12$ & \faTimes & \faTimes & \faTimes & & \faTimes \\ \hline
$-5$ & & \faTimes & \faTimes & & \faTimes \\ \hline
$\dfrac37$ & & & \faTimes & & \faTimes \\ \hline
$0,25$ & & & \faTimes & & \faTimes \\ \hline
$\sqrt{3}$ & & & & \faTimes & \faTimes \\ \hline
$\pi$ & & & & \faTimes & \faTimes \\ \hline
$-\sqrt{49}=-7$ & & \faTimes & \faTimes & & \faTimes \\ \hline
$0,\overline{3}$ & & & \faTimes & & \faTimes \\ \hline
$1,234\,567\,891\,011\ldots$ & & & & \faTimes & \faTimes \\ \hline
\end{tabular}\par}

\ex{2}
\begin{tasks}(5)
  \task $3\in\mathbb{N}$
  \task $-5\in\mathbb{Z}$
  \task $2\sqrt2\in\mathbb{I}$
  \task $\mathbb{Q}$
  \task $\mathbb{I}$
  \task $-3\in\mathbb{Z}$
  \task $0,3\in\mathbb{Q}$
  \task $\mathbb{I}$
  \task $-4\in\mathbb{Z}$
  \task $\mathbb{I}$
\end{tasks}

\ex{2}
\begin{tasks}(2)
  \task Vertadera: $\mathbb{N}\subset\mathbb{Z}$.
  \task Falsa: per exemple, $-3$.
  \task Vertadera: es pot escriure com a fracció.
  \task Falsa: per exemple, $\sqrt{4}=2$.
  \task Falsa: $\mathbb{Q}\cap\mathbb{I}=\varnothing$.
  \task Falsa: $\sqrt2=1,414\ldots$ és decimal i irracional.
\end{tasks}

\ex{3}
\begin{tasks}(1)
  \task Per exemple, $\pi$ o $\sqrt{10}\approx3,16$.
  \task Per exemple, $1,5$, ja que $1,5^{2}=2,25$ està entre 2 i 3.
  \task $n=1,\ 4,\ 9,\ 16,\ 25$ (els quadrats perfectes).
  \task No. Per exemple, $\sqrt2+\left(-\sqrt2\right)=0$, que és racional.
\end{tasks}

\end{document}
